% =============================================================================
% lusca-paper-latex default skeleton — generic article (English; pdflatex-ready)
%
% Usage: place your content between \begin{document} and \end{document}.
%   - For Chinese / CJK content: switch to xelatex and enable the xeCJK block
%     commented out below.
%   - To switch document class (IEEE / ACM / NeurIPS / apa7 ...): see
%     assets/templates/README.md.
%
% Compile:
%   latexmk -pdf main.tex
%   or four manual passes: pdflatex main && bibtex main && pdflatex main && pdflatex main
% =============================================================================

\documentclass[11pt,a4paper]{article}

% ---- Encoding & fonts (pdflatex) ----
\usepackage[utf8]{inputenc}
\usepackage[T1]{fontenc}
\usepackage{lmodern}
% For Chinese / CJK content: comment out the three lines above, enable the
% xeCJK block below, and compile with xelatex instead of pdflatex:
% \usepackage{xeCJK}
% \setCJKmainfont{Noto Serif CJK SC}
% \setCJKsansfont{Noto Sans CJK SC}
% \setCJKmonofont{Noto Sans Mono CJK SC}

% ---- Layout ----
\usepackage[margin=1in]{geometry}
\usepackage{microtype}

% ---- Math ----
\usepackage{amsmath}
\usepackage{amssymb}
\usepackage{amsthm}
\usepackage{mathtools}

% ---- Figures & tables ----
\usepackage{graphicx}
\usepackage{booktabs}      % three-line tables: \toprule \midrule \bottomrule
\usepackage{caption}
\usepackage{subcaption}    % subfigures
\usepackage{float}         % [H] for strict placement

% ---- Code (as needed) ----
\usepackage{listings}
\lstset{basicstyle=\ttfamily\small, breaklines=true, frame=single}

% ---- Bibliography ----
\usepackage[round,authoryear]{natbib}   % \citep \citet; use [square,numbers] for brackets

% ---- Links (hyperref MUST be loaded last) ----
\usepackage{xurl}           % break long URLs; load alongside/after hyperref
\usepackage{hyperref}
\hypersetup{colorlinks=true, linkcolor=blue, citecolor=blue, urlcolor=blue}

% ---- Theorem-like environments (enable as needed) ----
% \newtheorem{theorem}{Theorem}
% \newtheorem{lemma}{Lemma}
% \newtheorem{definition}{Definition}

% ---- Metadata ----
\title{Paper Title}
\author{%
  Author One\thanks{Corresponding author: \texttt{one@example.edu}} \and
  Author Two\\
  Department, Institution\\
  \texttt{\{one,two\}@example.edu}%
}
\date{\today}

\begin{document}
\maketitle

% Abstract: 150--250 words; present tense for general facts, past tense for what was done.
\begin{abstract}
Abstract paragraph (150--250 words): state the problem, the approach, the key results, and the main contribution.
\end{abstract}

\noindent\textbf{Keywords:} keyword one, keyword two, keyword three

% =============================================================================
\section{Introduction}
\label{sec:intro}

Cite with \citep{key2024} (parenthetical) or \citet{key2024} (narrative); cross-reference with Section~\ref{sec:method}, Figure~\ref{fig:arch}, Table~\ref{tab:results}, Equation~\eqref{eq:loss}. Use a non-breaking space (\textasciitilde{}) to keep a number attached to its label.

% =============================================================================
\section{Method}
\label{sec:method}

Inline math such as $\mathcal{L}(\theta)$ is delimited by dollar signs. Display equations use the \texttt{equation} environment with a \textbackslash label for cross-referencing:
\begin{equation}
  \label{eq:loss}
  \mathcal{L}(\theta) = -\frac{1}{N}\sum_{i=1}^{N} \log p_\theta(y_i \mid x_i).
\end{equation}

Multi-line alignments use \texttt{align}:
\begin{align}
  \nabla_\theta \mathcal{L} &= \mathbb{E}_{(x,y)\sim\mathcal{D}}\bigl[\nabla_\theta \log p_\theta(y\mid x)\bigr], \label{eq:grad}\\
  \theta &\leftarrow \theta - \eta\, \nabla_\theta \mathcal{L}. \label{eq:update}
\end{align}

% ---- Figure example (uncomment as needed) ----
% \begin{figure}[htbp]
%   \centering
%   \includegraphics[width=0.8\linewidth]{figures/arch}
%   \caption{Overview of the architecture. Write captions as full declarative sentences, capitalized, with a closing period.}
%   \label{fig:arch}
% \end{figure}

% ---- Table example (uncomment as needed) ----
% \begin{table}[htbp]
%   \centering
%   \caption{Results on the benchmark.}
%   \label{tab:results}
%   \begin{tabular}{lcc}
%     \toprule
%     Method & Accuracy & F1 \\
%     \midrule
%     Baseline & 0.72 & 0.70 \\
%     Ours     & \textbf{0.89} & \textbf{0.87} \\
%     \bottomrule
%   \end{tabular}
% \end{table}

% =============================================================================
\section{Experiments}
\label{sec:exp}

Report results in the past tense (``The model achieved \ldots''), but refer to tables and figures in the present tense (``Table~\ref{tab:results} shows \ldots'').

% =============================================================================
\section{Conclusion}
\label{sec:conclusion}

Summarize the contribution and its significance. Do not merely repeat the abstract.

% =============================================================================
% Bibliography --- choose ONE option. The default is the inline Option B so the
% template compiles out-of-the-box; switch to Option A when you have a .bib file.

% --- Option A: BibTeX (recommended when references.bib is present) ---
% \bibliographystyle{plainnat}
% \bibliography{references}

% --- Option B: inline (self-contained; compiles without bibtex) ---
\begin{thebibliography}{99}
\bibitem[Author(2024)]{key2024}
  Author, A., \& Coauthor, B. (2024). Paper title. \emph{Journal Name}, 12(3), 45--67.
\end{thebibliography}

% =============================================================================
% Typeset with lusca-paper-latex (LaTeX comments; safe to delete before submission).
% Author: lusca
% Version: lusca-paper-latex v1.0.0
% Source: https://github.com/yjmm10/lusca-skill/tree/main/skills/lusca-paper-latex
% =============================================================================

\end{document}
